\documentclass[11pt]{article}
\usepackage[a4paper,margin=1in]{geometry}
\usepackage{amsmath,amssymb,amsthm,booktabs}
\usepackage[T1]{fontenc}
\usepackage{lmodern}
\usepackage[hidelinks]{hyperref}
\newtheorem{theorem}{Theorem}
\newtheorem{proposition}{Proposition}
\newcommand{\Z}{\mathbb Z}
\newcommand{\Fix}{\operatorname{Fix}}
\newcommand{\FC}{\operatorname{FC}}
\newcommand{\Sset}{\mathcal S}
\title{A counterexample to ideal closure of property-(s) elements in skew braces}
\author{Alec Kriebel\\
\small Independent researcher\\
\small \href{https://orcid.org/0009-0001-9320-500X}{ORCID: 0009-0001-9320-500X}}
\date{12 October 2026}
\begin{document}
\maketitle
\begin{abstract}
The elements with property (s) in a skew left brace need not form an
additive subgroup, even when the additive group is abelian.
On $\Z\times\Z/4\Z$ we give an explicit standard semidirect brace whose
property-(s) set is $\Z\times\{0,1\}$. The element $(0,1)$ has both
defining indices equal to two, whereas its additive double has both
indices infinite. This answers the unrestricted ideal-closure question
of Colazzo negatively. The construction uses known ingredients.
We credit an exact overlapping October 2026 preprint and record our
earlier September 2026 repository disclosure; historical priority and
mutual independence remain unresolved.
\end{abstract}

\section{The question and exact conventions}
A skew left brace is a set $B$ with group operations $+$ and $\circ$
satisfying
\[
 a\circ(b+c)=(a\circ b)-a+(a\circ c).
\]
The common identity is denoted by $0$; addition need not be commutative.
Write $\lambda_a(b)=-a+a\circ b$ and $a*b=\lambda_a(b)-b$.
For $z\in B$, put
\[
 H_z=\{g:z*g=0\}\cap C_{(B,+)}(z),\qquad
 K_z=\{g:g*z=0\}\cap C_{(B,\circ)}(z).
\]
An element has \emph{property (s)} when both
\begin{equation}\label{indices}
 [(B,+):H_z]<\infty,\qquad [(B,\circ):K_z]<\infty .
\end{equation}
These are the right and left fixed intersections in the original
definition; their ambient groups are different.

Colazzo's Question~1 in Oberwolfach Report 9/2023, printed
p.~543~\cite{OWR}, asks whether the set $\Sset(B)$ of these elements is
an ideal for an arbitrary skew brace.
The question imposes no finiteness, global property-(S), finite-generation
or two-sided assumption. In a finite brace every element has (s).
Colazzo--Ferrara--Trombetti's published Definition~3.7 uses
\eqref{indices}; Proposition~3.10 proves ideality for two-sided braces
and the following text leaves the general case open~\cite{CFT}.
An ideal must in particular be an additive subgroup.

\section{The counterexample and complete index calculation}
On $C=\Z/4\Z$ set
\[
 c\diamond d=c+d+2cd\pmod4,\qquad
 \varepsilon(0)=\varepsilon(1)=1,\quad
 \varepsilon(2)=\varepsilon(3)=-1.
\]
Residues in statements involving $\{0,1\}$ are represented by
$0,1,2,3$.

\begin{theorem}\label{main}
Give $B=\Z\times C$ coordinatewise addition and multiplication
\begin{equation}\label{circle}
 (n,c)\circ(m,d)=
 \bigl(n+\varepsilon(c)m,\ c\diamond d\bigr).
\end{equation}
Then $B$ is a left brace with abelian additive group, and
\[
 \Sset(B)=\Z\times\{0,1\}.
\]
For $x=(0,1)$ both indices in \eqref{indices} equal $2$.
For $x+x=(0,2)$ both are infinite.
Thus $\Sset(B)$ is neither an additive subgroup nor an ideal.
It also fails $\lambda$-invariance, although it is a normal subgroup of
$(B,\circ)$.
\end{theorem}

\begin{proof}
Write $c=u+2v$, $d=s+2t$, with all four bits in $\{0,1\}$.
Modulo four,
\[
 c\diamond d=(u+s\bmod2)+2(v+t\bmod2).
\]
Indeed the carry from $u+s$ is canceled modulo four by $2us$.
Hence $(C,\diamond)$ is the Klein four group and
$\varepsilon(c)=(-1)^v$ is a character of this group.
Equation \eqref{circle} is its action by signs on $\Z$, so it is an
associative group law with identity $(0,0)$ and inverse
\[
 (n,c)^{-1}_{\circ}=(-\varepsilon(c)n,c).
\]
In particular, this is a group construction over all integers.

Direct calculation gives
\begin{align}
 \lambda_{(n,c)}(m,d)&=(\varepsilon(c)m,(-1)^c d),\label{lambda}\\
 (n,c)*(m,d)&=((\varepsilon(c)-1)m,2cd).\label{star}
\end{align}
The second coordinates are modulo four.
Each map in \eqref{lambda} is an additive automorphism. Together
with $a\circ b=a+\lambda_a(b)$ this proves the left brace identity.

For $z=(n,c)$ and $g=(m,d)$, addition is abelian, so
\begin{align}
 H_z&=\{(m,d):(\varepsilon(c)-1)m=0,\quad 2cd=0\},\label{H}\\
 K_z&=\{(m,d):(\varepsilon(c)-1)m=0,\quad
             (\varepsilon(d)-1)n=0,\quad 2cd=0\}.\label{K}
\end{align}
To verify \eqref{K}, $g*z=0$ imposes
$(\varepsilon(d)-1)n=0$ and $2cd=0$.
The remaining commutation equation is
$(\varepsilon(c)-1)m=(\varepsilon(d)-1)n$,
which then reduces to the first equation in \eqref{K}.

For $c=0,1$, these groups have the form $\Z\times L$.
Their exact indices are
\[
\begin{array}{c|cc}
 z & [(B,+):H_z] & [(B,\circ):K_z]\\ \hline
 (0,0)&1&1\\
 (n,0),\ n\ne0&1&2\\
 (0,1)&2&2\\
 (n,1),\ n\ne0&2&4 .
\end{array}
\]
For the second index the respective subsets $L$ are
$C,\{0,1\},\{0,2\},\{0\}$, all subgroups of $(C,\diamond)$.
The projection $(B,\circ)\to(C,\diamond)$ is surjective, so these
indices are $4/|L|$.
For the first index, \eqref{H} gives $B$ when $c=0$ and
$\Z\times\{0,2\}$ when $c=1$.

If $c=2$ or $3$, both \eqref{H} and \eqref{K} force $m=0$.
The additive cosets represented by $(j,0)$, $j\in\Z$, are pairwise
distinct, since their differences have nonzero first coordinate.
The multiplicative cosets are also pairwise distinct:
$(i,0)^{-1}_{\circ}\circ(j,0)=(j-i,0)$ belongs to $K_z$ only when
$i=j$. This proves both infinite indices and the asserted complete
set, without extrapolation from finite computations.

Finally $x+x=(0,2)$, while
$\lambda_x(x)=(0,3)\notin\Sset(B)$.
The character $(n,c)\mapsto\varepsilon(c)$ on $(B,\circ)$ has
kernel $\Sset(B)$, proving multiplicative normality.
\end{proof}

The circle group is $D_\infty\times C_2$: the high bit acts by
inversion on $\Z$, while the low bit acts trivially.
The brace is not two-sided. For $u=v=(0,1)$ and $w=(1,0)$,
\[
 (u+v)\circ w=(-1,2),\qquad
 (u\circ w)-w+(v\circ w)=(1,2).
\]
Thus the counterexample respects the two-sided restriction in the
known positive result.

\section{Known ingredients, overlap and priority limits}
The finite brace $(C,+,\diamond)$ is old:
Bachiller's Proposition~2.4(iii) includes it~\cite{Bachiller}, and
CFT Example~3.8 recalls it~\cite{CFT}.
The semidirect construction is standard, for example
Smoktunowicz--Vendramin's published Corollary~3.37~\cite{SV}.
Dietzel's Theorem~5(iv) explicitly uses the same high-bit inversion
character in a finite cyclic-module family~\cite{Dietzel}.
Replacing that trivial finite module by the trivial brace $\Z$
and making the index calculation above is an elementary specialization.
We claim no new construction method.

The example is also an immediate application of the earlier
finite-conjugacy machinery once the character is chosen.
Cascella--Properzi--Van Antwerpen's Proposition~5.17 identifies
$\Sset(B)$ with the multiplicative FC-centre when every element has
finite additive conjugacy and a finite $\lambda$-orbit~\cite{CPV}.
Those hypotheses hold here: addition is abelian and the $\lambda$
image has four elements. Conjugation is explicitly
\[
 (m,d)\circ(n,c)\circ(m,d)^{-1}_{\circ}
   =\bigl((1-\varepsilon(c))m+\varepsilon(d)n,c\bigr),
\]
which gives the same set. The direct proof of
Theorem~\ref{main} does not depend on that characterization.

Di Matteo--Ferrara also give a pre-September characterization and
a finite-factor construction~\cite{DMF}.
An infinite-module variation of their Section~4, and another of
Mondal--Shumyatsky--Trombetti--Yadav's finite Theorem~4.7~\cite{MSTY},
give alternative negative examples by short calculations.
The supporting priority audit records those derivations.
They are deductions made in our audit, not claims that those papers
explicitly published the unrestricted negative answer.
This ancestry limits any broad originality claim.

Alper Ferudun's unrefereed October~11, 2026 note,
Theorem~1.2~\cite{Ferudun}, gives exactly the same construction and
property-(s) set as Theorem~\ref{main}. We credit this exact overlap.
GitHub records our PR~156 as created on September~30, 2026 at
12:24:37 UTC~\cite{Repository}; its original proof is unchanged.
The first proof commit is dated September~30 at 12:16:29 UTC, an
author-supplied commit timestamp distinct from the provider's PR date.
Ferudun's first Zenodo record was created on October~11 at
05:15:40.640397 UTC. This note's formal preprint date is later and is
not backdated. The repository was public according to the author;
current public visibility was checked, while historical visibility
was not independently archived.

We describe these as \emph{contemporaneous separate discoveries},
meaning separate disclosures in the same period. The dates establish
an ordering of the recorded disclosures, not a literally simultaneous
release or either author's awareness of the other's work.
Mutual independence has not been established.

Our bounded primary-source audit found no earlier explicit negative
answer in the full bodies read. It does not certify historical novelty,
first priority or exclusive resolution.
Several complete publisher bodies were unavailable: JKVV's 2021
\emph{Radical and weight} article, the 2026 CPV and Mondal--Yadav
articles, and the 2026 word/conjugacy-problem article
(DOI \href{https://doi.org/10.1080/10586458.2025.2598349}
{10.1080/10586458.2025.2598349}).
Available full preprints or the author-labeled CPV postprint were
read where available. The precise versions, additional access gaps
and bounded search limits are preserved in the supporting audit.
These gaps can affect priority assessment; they do not affect the
self-contained counterexample.

\section*{Verification and disclosure}
The support archive contains the standalone source, an exact
standard-library checker, recorded ordinary and optimized runs,
negative controls, source locators and the priority/provenance audit.
The checker uses explicit raised errors, so optimization does not
disable its checks. Its finite and symbolic controls supplement the
written proof; the infinite-coset argument is the certificate of
infinite index. Every finite quotient has property (s) globally and
cannot by itself establish this failure.

AI tools were used extensively in solving, drafting, literature
comparison, reproduction and adversarial verification.
This preprint is unrefereed and has not undergone conventional human
peer review. The author is responsible for its contents.

\begin{thebibliography}{10}
\bibitem{OWR}
I.~Colazzo, Derived-indecomposable solutions and skew braces whose
elements have a finite number of conjugates, contribution to
\emph{Oberwolfach Report} 9/2023, pp.~543--544, Question~1.
\href{https://doi.org/10.4171/OWR/2023/9}{doi:10.4171/OWR/2023/9}.

\bibitem{CFT}
I.~Colazzo, M.~Ferrara and M.~Trombetti,
On derived-indecomposable solutions of the Yang--Baxter equation,
\emph{Publ. Mat.} 69 (2025), 171--193.
\href{https://doi.org/10.5565/PUBLMAT6912508}{doi:10.5565/PUBLMAT6912508}.

\bibitem{Bachiller}
D.~Bachiller, Classification of braces of order $p^3$,
\emph{J. Pure Appl. Algebra} 219 (2015), 3568--3603.
\href{https://doi.org/10.1016/j.jpaa.2014.12.013}
{doi:10.1016/j.jpaa.2014.12.013};
\href{https://arxiv.org/abs/1407.5224v1}{arXiv:1407.5224v1}.

\bibitem{SV}
A.~Smoktunowicz and L.~Vendramin,
On skew braces (with an appendix by N.~Byott and L.~Vendramin),
\emph{J. Comb. Algebra} 2 (2018), 47--86.
\href{https://doi.org/10.4171/JCA/2-1-3}{doi:10.4171/JCA/2-1-3}.

\bibitem{Dietzel}
C.~Dietzel, Braces of order $p^2q$,
\emph{J. Algebra Appl.} 20 (2021), 2150140.
\href{https://doi.org/10.1142/S0219498821501401}
{doi:10.1142/S0219498821501401};
\href{https://arxiv.org/abs/1801.06911v2}{arXiv:1801.06911v2},
Theorem~5(iv).

\bibitem{CPV}
R.~Cascella, S.~Properzi and A.~Van Antwerpen,
Finiteness conditions on skew braces and solutions of the
Yang--Baxter equation, \emph{Mediterr. J. Math.} 23 (2026), article~187.
\href{https://doi.org/10.1007/s00009-026-03182-4}
{doi:10.1007/s00009-026-03182-4}.
Full author-labeled postprint:
\href{https://arxiv.org/abs/2603.06177v2}{arXiv:2603.06177v2},
submitted September~14, 2026.

\bibitem{DMF}
M.~Di Matteo and M.~Ferrara, Finite-index problems in skew braces,
\href{https://arxiv.org/abs/2607.20222v1}{arXiv:2607.20222v1},
July~22, 2026.

\bibitem{MSTY}
S.~Mondal, P.~Shumyatsky, M.~Trombetti and M.~K.~Yadav,
Relative commuting probability and BFC-type results for finite
skew braces,
\href{https://arxiv.org/abs/2609.16269v1}{arXiv:2609.16269v1},
September~14, 2026.

\bibitem{Ferudun}
A.~Ferudun, The Elements with Property (s) of a Skew Brace Need Not
Form an Ideal: On a Question of Colazzo, Ferrara and Trombetti,
unrefereed preprint, version~1, October~11, 2026.
\href{https://doi.org/10.5281/zenodo.23292254}
{doi:10.5281/zenodo.23292254}.

\bibitem{Repository}
A.~Kriebel, Original counterexample draft, PR~156,
\href{https://github.com/AlecKriebel/Math/pull/156}
{github.com/AlecKriebel/Math/pull/156}, September~30, 2026.
Initial proof commit
\href{https://github.com/AlecKriebel/Math/commit/c1677fe1ed901b8d9c5efe4cbd3d3fec7d0ba227}
{\texttt{c1677fe1ed901b8d9c5efe4cbd3d3fec7d0ba227}}.
\end{thebibliography}
\end{document}

